% GENERATED FILE. Edit canonical lessons, then run npm run generate:books.
% canonical-source-hash: fnv1a64:6d63b9cdba52e5f1
% canonical-lessons: GE-C372-vereinbaren, GE-C372-vereinbarung, GE-C372-ausmachen, GE-C372-zeitplan, GE-C372-umplanen

\chapter{Etwas vereinbaren}
\label{ch:ge-etwas-vereinbaren}
\hlchaptermodality{372}

\begin{chapteropening}
\textbf{By the end of this chapter:} \emph{I can agree on a date or an arrangement, keep to a schedule and change my plans.}

Complete the last lesson of chapter 372: I can agree on a date or an arrangement, keep to a schedule and change my plans.
\end{chapteropening}

\section[vereinbaren]{vereinbaren — to agree on, to arrange (formally)}

\label{lesson:GE-C372-vereinbaren}



\begin{quote}
Before the new one: say the German for the rejection, then the German for the confirmation.
\end{quote}

\subsection*{You'll want to know: vereinbaren}
\textbf{vereinbaren} — \textquotedblleft{}to agree on, to arrange (formally)\textquotedblright{}.

\begin{quote}
Wir haben einen Termin für Dienstag vereinbart.
\end{quote}

\begin{quote}
\emph{We have arranged an appointment for Tuesday.}
\end{quote}

Inseparable: \emph{wir vereinbaren, wir haben vereinbart}.

\begin{cousinweb}
\emph{vereinbaren} is inherited from Middle High German \emph{vereinbæren} (EtymDB), built on \emph{einbære} \textquotedblleft{}united\textquotedblright{}, from \emph{ein} \textquotedblleft{}one\textquotedblright{}, Proto-Germanic \emph{ainaz}: two sides become one.
\end{cousinweb}

\begin{grammarlens}[title={What you've built}]
One more word for telling, explaining and describing at B1.
\end{grammarlens}

\subsection*{Your turn}
\begin{itemize}
\raggedright
  \item \emph{Say it:} \emph{vereinbaren}
  \item \emph{Say it:} \emph{vereinbaren}, once more
  \item \emph{Say it:} a sentence of your own with \emph{vereinbaren}
  \item \emph{Recall:} say the German for the rejection, then the German for the confirmation, then say \emph{vereinbaren} again
\end{itemize}

\begin{tcolorbox}[breakable,colback=teal!4,colframe=teal!35!black,title={Before you move on}]
What does vereinbaren mean? (\textquotedblleft{}To agree on, to arrange (formally)\textquotedblright{}.) Say it once more.
\end{tcolorbox}
\section[die Vereinbarung]{die Vereinbarung — the arrangement, the agreement}

\label{lesson:GE-C372-vereinbarung}



\begin{quote}
Before the new one: say the German for the confirmation, then the German for to agree on.
\end{quote}

\subsection*{You'll want to know: die Vereinbarung}
\textbf{die Vereinbarung} — \textquotedblleft{}the arrangement, the agreement\textquotedblright{}.

\begin{quote}
Nach unserer Vereinbarung bezahlt jeder die Hälfte.
\end{quote}

\begin{quote}
\emph{According to our arrangement, each of us pays half.}
\end{quote}

Plural: \emph{die Vereinbarungen}. \emph{Nach Vereinbarung} — by appointment.

\begin{cousinweb}
\emph{Vereinbarung} is the noun of \emph{vereinbaren}, from Middle High German \emph{vereinbæren} (EtymDB), built on \emph{ein} \textquotedblleft{}one\textquotedblright{}.
\end{cousinweb}

\begin{grammarlens}[title={What you've built}]
One more word for telling, explaining and describing at B1.
\end{grammarlens}

\subsection*{Your turn}
\begin{itemize}
\raggedright
  \item \emph{Say it:} \emph{die Vereinbarung}
  \item \emph{Say it:} \emph{die Vereinbarung}, once more
  \item \emph{Say it:} a sentence of your own with \emph{die Vereinbarung}
  \item \emph{Recall:} say the German for the confirmation, then the German for to agree on, then say \emph{die Vereinbarung} again
\end{itemize}

\begin{tcolorbox}[breakable,colback=teal!4,colframe=teal!35!black,title={Before you move on}]
What does die Vereinbarung mean? (\textquotedblleft{}The arrangement, the agreement\textquotedblright{}.) Say it once more.
\end{tcolorbox}
\section[ausmachen]{ausmachen — to arrange (informal); to switch off; to matter}

\label{lesson:GE-C372-ausmachen}



\begin{quote}
Before the new one: say the German for to agree on, then the German for the arrangement.
\end{quote}

\subsection*{You'll want to know: ausmachen}
\textbf{ausmachen} — \textquotedblleft{}to arrange (informal); to switch off; to matter\textquotedblright{}.

\begin{quote}
Wir haben ausgemacht, dass wir uns um acht treffen.
\end{quote}

\begin{quote}
\emph{We arranged to meet at eight.}
\end{quote}

Separable: \emph{wir machen aus}, \emph{wir haben ausgemacht}. Also to switch off: \emph{das Licht ausmachen}; and to matter: \emph{Das macht nichts aus} — that does not matter.

\begin{cousinweb}
\emph{ausmachen} is built from \emph{aus} \textquotedblleft{}out\textquotedblright{} and \emph{machen} \textquotedblleft{}to make\textquotedblright{} (EtymDB): you make a plan out between you.
\end{cousinweb}

\begin{grammarlens}[title={What you've built}]
One more word for telling, explaining and describing at B1.
\end{grammarlens}

\subsection*{Your turn}
\begin{itemize}
\raggedright
  \item \emph{Say it:} \emph{ausmachen}
  \item \emph{Say it:} \emph{ausmachen}, once more
  \item \emph{Say it:} a sentence of your own with \emph{ausmachen}
  \item \emph{Recall:} say the German for to agree on, then the German for the arrangement, then say \emph{ausmachen} again
\end{itemize}

\begin{tcolorbox}[breakable,colback=teal!4,colframe=teal!35!black,title={Before you move on}]
What does ausmachen mean? (\textquotedblleft{}To arrange (informal); to switch off; to matter\textquotedblright{}.) Say it once more.
\end{tcolorbox}
\section[der Zeitplan]{der Zeitplan — the schedule}

\label{lesson:GE-C372-zeitplan}



\begin{quote}
Before the new one: say the German for the arrangement, then the German for to arrange.
\end{quote}

\subsection*{You'll want to know: der Zeitplan}
\textbf{der Zeitplan} — \textquotedblleft{}the schedule\textquotedblright{}.

\begin{quote}
Wir sind gut im Zeitplan: Das Haus ist im Juni fertig.
\end{quote}

\begin{quote}
\emph{We are well on schedule: the house will be finished in June.}
\end{quote}

Plural: \emph{die Zeitpläne}. For trains it is \emph{der Fahrplan}.

\begin{cousinweb}
\emph{Zeitplan} joins \emph{Zeit} \textquotedblleft{}time\textquotedblright{}, Proto-Germanic \emph{tīdiz} (EtymDB), to \emph{Plan}, borrowed from French \emph{plan} and going back to Latin \emph{plānus} \textquotedblleft{}level, flat\textquotedblright{} (EtymDB; Lewis \& Short): a drawing laid out flat.
\end{cousinweb}

\begin{grammarlens}[title={What you've built}]
One more word for telling, explaining and describing at B1.
\end{grammarlens}

\subsection*{Your turn}
\begin{itemize}
\raggedright
  \item \emph{Say it:} \emph{der Zeitplan}
  \item \emph{Say it:} \emph{der Zeitplan}, once more
  \item \emph{Say it:} a sentence of your own with \emph{der Zeitplan}
  \item \emph{Recall:} say the German for the arrangement, then the German for to arrange, then say \emph{der Zeitplan} again
\end{itemize}

\begin{tcolorbox}[breakable,colback=teal!4,colframe=teal!35!black,title={Before you move on}]
What does der Zeitplan mean? (\textquotedblleft{}The schedule\textquotedblright{}.) Say it once more.
\end{tcolorbox}
\section[umplanen]{umplanen — to change your plans}

\label{lesson:GE-C372-umplanen}



\begin{quote}
Before the new one: say the German for to arrange, then the German for the schedule.
\end{quote}

\subsection*{You'll want to know: umplanen}
\textbf{umplanen} — \textquotedblleft{}to change your plans\textquotedblright{}.

\begin{quote}
Wegen des Regens mussten wir umplanen.
\end{quote}

\begin{quote}
\emph{Because of the rain we had to change our plans.}
\end{quote}

Separable: \emph{wir planen um}, \emph{wir haben umgeplant}.

\begin{cousinweb}
\emph{umplanen} is built from \emph{um} \textquotedblleft{}around, anew\textquotedblright{} and \emph{planen} \textquotedblleft{}to plan\textquotedblright{}, from \emph{Plan}, borrowed from French \emph{plan} and going back to Latin \emph{plānus} \textquotedblleft{}flat\textquotedblright{} (EtymDB).
\end{cousinweb}

\begin{grammarlens}[title={What you've built}]
One more word for telling, explaining and describing at B1.
\end{grammarlens}

\subsection*{Your turn}
\begin{itemize}
\raggedright
  \item \emph{Say it:} \emph{umplanen}
  \item \emph{Say it:} \emph{umplanen}, once more
  \item \emph{Say it:} a sentence of your own with \emph{umplanen}
  \item \emph{Recall:} say the German for to arrange, then the German for the schedule, then say \emph{umplanen} again
\end{itemize}

\begin{tcolorbox}[breakable,colback=teal!4,colframe=teal!35!black,title={Before you move on}]
What does umplanen mean? (\textquotedblleft{}To change your plans\textquotedblright{}.) Say it once more.
\end{tcolorbox}
